\section{The Insight after the Scholastics}

The recognition did not disappear when medieval vocabulary ceased to govern philosophy. In the Third Meditation, Descartes argues that past existence does not by itself account for present existence. Preservation and creation differ only conceptually in the respect under discussion: a cause is required not only for the self's production but for its conservation now. His larger method and metaphysics differ sharply from Aquinas's, but the present tense of dependence survives.

The Catholic doctrinal boundary was stated with particular force in the nineteenth century. Vatican I confesses a God really and essentially distinct from the world, freely creating not to acquire or increase divine beatitude but to manifest perfection through goods given to creatures (\emph{Dei Filius}, ch.~1). When the constitution turns to faith, it says that the human person wholly depends on God as Creator and Lord (ch.~3). The council defines the Creator--creature distinction and free creation; it does not canonize every Thomistic explanation used in this paper.

The \emph{Catechism of the Catholic Church} gives the contemporary synthesis. God is at once infinitely greater than his works and present to their inmost being (CCC 300). He sustains them at every moment, enables their action, and orders them to their end (CCC 301). This ``utter dependence'' is named a source not only of wisdom but of freedom, joy, and confidence. The next articles safeguard secondary causality and human cooperation.

Across these changes of idiom, the constant is not one image or proof. It is a structure: finite actuality does not explain itself; dependence reaches the present; the source is not reciprocally dependent; and the received act is sufficiently real to ground genuine action. The phrase \emph{asymmetric dependence} is contemporary. The structure it names is old.
